% Clase 28 — curvas de resonancia del RLC serie para tres resistencias, en
% unidades de omega0 = 1/sqrt(LC). Con x = omega/omega0 y Q = omega0 L / R:
%   amplitud  I0/(V0/R) = 1/sqrt(1 + Q^2 (x - 1/x)^2)   ... (normalizada al
%             máximo de cada curva; el pico absoluto V0/R crece al bajar R)
%   fase      theta = arctan(Q (x - 1/x))
\begin{tikzpicture}
  \begin{axis}[
    emfig, width=6.9cm, height=5.4cm,
    xmin=0.3, xmax=2, ymin=0, ymax=1.1,
    xtick={0.5,1,1.5,2}, ytick={0,0.5,0.707,1},
    yticklabels={$0$,$0{,}5$,$1/\sqrt2$,$1$},
    xlabel={$\omega/\omega_0$}, ylabel={$I_0/I_0(\omega_0)$},
    title={\small (a) amplitud},
    legend style={at={(1.1,-0.42)}, anchor=north, legend columns=3,
      /tikz/every even column/.append style={column sep=1em}},
  ]
    \addplot[figblue!40, domain=0.3:2, samples=200] {1/sqrt(1+4*(x-1/x)^2)};
    \addlegendentry{$\mathcal{Q} = 2$}
    \addplot[figblue!75, domain=0.3:2, samples=200] {1/sqrt(1+25*(x-1/x)^2)};
    \addlegendentry{$\mathcal{Q} = 5$}
    \addplot[figblue, domain=0.3:2, samples=400] {1/sqrt(1+400*(x-1/x)^2)};
    \addlegendentry{$\mathcal{Q} = 20$}
  \end{axis}
  \begin{axis}[
    emfig, width=6.9cm, height=5.4cm, xshift=7.6cm,
    xmin=0.3, xmax=2, ymin=-1.7, ymax=1.7,
    xtick={0.5,1,1.5,2}, ytick={-1.5708,0,1.5708},
    yticklabels={$-\pi/2$,$0$,$\pi/2$},
    xlabel={$\omega/\omega_0$}, ylabel={$\theta$},
    title={\small (b) fase},
  ]
    \addplot[figblue!40, domain=0.3:2, samples=200] {rad(atan(2*(x-1/x)))};
    \addplot[figblue!75, domain=0.3:2, samples=200] {rad(atan(5*(x-1/x)))};
    \addplot[figblue, domain=0.3:2, samples=400] {rad(atan(20*(x-1/x)))};
  \end{axis}
\end{tikzpicture}
